% !TEX TS-program = pdflatex
% !TeX program = pdflatex
% !TEX encoding = UTF-8
% !TEX spellcheck = en_US

\documentclass[xcolor=table]{beamer}

\usepackage{../extra/beamer/karimnlp}

\input{options}

\subtitle[07- Word semantics]{Chapter 07\\Word semantics} 

\changegraphpath{../img/word-sem/}

\begin{document}
	
\begin{frame}
	\frametitle{\inserttitle}
	\framesubtitle{\insertshortsubtitle: Introduction}

	\begin{exampleblock}{Example of polysemy}
		\begin{center}
			\Large\bfseries
	%		\begin{itemize}
				The hotel has a big \expword{wing}.
				 
				That bird broke its \expword{wing}.
				
				They attacked the enemy's right \expword{wing}.
				
				I ate just a \expword{wing}.
	%		\end{itemize}
		\end{center}
	\end{exampleblock}
	
	\begin{itemize}
		\item Do all occurrences of the word ``\expword{wing}'' have the same meaning?
		\item How can we determine the intended meaning of a word in context?
		\item Possible senses of ``\expword{wing}'': \url{http://wordnetweb.princeton.edu/perl/webwn?s=wing}
	\end{itemize}

\end{frame}

\begin{frame}
	\frametitle{\inserttitle}
	\framesubtitle{\insertshortsubtitle: Plan}

	\begin{multicols}{2}
	%	\small
	\tableofcontents
	\end{multicols}

\end{frame}

%===================================================================================
\section{Lexical databases}
%===================================================================================

\begin{frame}
	\frametitle{\insertshortsubtitle}
	\framesubtitle{\insertsection}
	
	\begin{itemize}
		\item A lexical database is a structured resource that stores a language's vocabulary along with linguistic information.
		\item Lexical information may include:
		\begin{itemize}
			\item PoS tags, lemmas, frequencies;
			\item Morphological and syntactic properties;
			\item Semantic relations (e.g., \expword{synonymy, hypernymy}).
		\end{itemize}
	\end{itemize}
	
	\begin{figure}
		\centering 
		\hgraphpage[.5\textwidth]{exp-bd-lex.pdf}
		\caption{Example of categories and their relations in a lexical database \cite{2019-white-al}}
	\end{figure}
	
\end{frame}


\subsection{Semantic relations}

\begin{frame}
	\frametitle{\insertshortsubtitle: \insertsection}
	\framesubtitle{\insertsubsection: Reminder}
	
	\begin{itemize}
		\item \optword{Synonymy}: words that have similar meanings (in a given context).
		\item \optword{Antonymy}: words that have opposite meanings (in a given context).
		\item Taxonomic (classification) relations:
		\begin{itemize}
			\item \optword{Hyponymy}: a relation where a concept is more specific than another. It corresponds to an \keyword{IS-A} relation. E.g. \expword{car IS-A vehicle}.
			\item \optword{Hypernymy}: a relation where a concept is more general than another.
			\item \optword{Meronymy}: a part–whole relation. E.g. \expword{wheel is part of a car; car is the holonym of wheel}.
		\end{itemize}
	\end{itemize}
	
\end{frame}


\subsection{WordNet}

\begin{frame}
	\frametitle{\insertshortsubtitle: \insertsection}
	\framesubtitle{\insertsubsection}

	\vspace{-6pt}
	\begin{exampleblock}{Example from WordNet: \small\url{http://wordnetweb.princeton.edu/perl/webwn}}
		\fontsize{6}{6}\selectfont
		%		\begin{alltt}
		{\small\bfseries Noun}
		\begin{itemize}\setlength\itemsep{0pt}
			\item \textcolor{blue}{\underline{S:}} \textcolor{red}{(n)} \textbf{reason}, \textcolor{blue}{\underline{ground}} (a rational motive for a belief or action) \textit{"the reason that war was declared"; "the grounds for their declaration"}
			\item \textcolor{blue}{\underline{S:}} \textcolor{red}{(n)} \textbf{reason} (an explanation of the cause of some phenomenon) \textit{"the reason a steady state was never reached was that the back pressure built up too slowly"}
			\item \textcolor{blue}{\underline{S:}} \textcolor{red}{(n)} \textbf{reason}, \textcolor{blue}{\underline{understanding}}, \textcolor{blue}{\underline{intellect}} (the capacity for rational thought or inference or discrimination) \textit{"we are told that man is endowed with reason and capable of distinguishing good from evil"}
			\item \textcolor{blue}{\underline{S:}} \textcolor{red}{(n)} \textcolor{blue}{\underline{rationality}}, \textbf{reason}, \textcolor{blue}{\underline{reasonableness}} (the state of having good sense and sound judgment) \textit{"his rationality may have been impaired"; "he had to rely less on reason than on rousing their emotions"}
			\item \textcolor{blue}{\underline{S:}} \textcolor{red}{(n)} \textcolor{blue}{\underline{cause}}, \textbf{reason}, \textcolor{blue}{\underline{grounds}} (a justification for something existing or happening) \textit{"he had no cause to complain"; "they had good reason to rejoice"}
			\item \textcolor{blue}{\underline{S:}} \textcolor{red}{(n)} \textbf{reason} (a fact that logically justifies some premise or conclusion) \textit{"there is reason to believe he is lying"}
		\end{itemize}
		
		{\small\bfseries Verb}
		\begin{itemize}\setlength\itemsep{0pt}
			\item \textcolor{blue}{\underline{S:}} \textcolor{red}{(v)} \textbf{reason}, \textcolor{blue}{\underline{reason out}}, \textcolor{blue}{\underline{conclude}} (decide by reasoning; draw or come to a conclusion) \textit{"We reasoned that it was cheaper to rent than to buy a house"}
			\item \textcolor{blue}{\underline{S:}} \textcolor{red}{(v)} \textcolor{blue}{\underline{argue}}, \textbf{reason} (present reasons and arguments)
			\item \textcolor{blue}{\underline{S:}} \textcolor{red}{(v)} \textbf{reason} (think logically) \textit{"The children must learn to reason"}
		\end{itemize}
	\end{exampleblock}
	
\end{frame}

\begin{frame}
	\frametitle{\insertshortsubtitle: \insertsection}
	\framesubtitle{\insertsubsection: Description}
	
	\begin{itemize}
		\item A lexical database for English \cite{1995-miller}
		\item Four parts of speech: nouns, verbs, adjectives, and adverbs
		\item Each sense is represented by a synset identifier (\keyword{synset}).
		\begin{itemize}
			\item \keyword{synset} (set of synonyms): groups words that share the same sense.
			\item E.g. \expword{05659525: reason\#3, understanding\#4, intellect\#2}
		\end{itemize}
		\item Each sense is defined by a \keyword{gloss}.
		\begin{itemize}
			\item E.g. \expword{05659525: (the capacity for rational thought or inference or discrimination) "we are told that man is endowed with reason and capable of distinguishing good from evil"}.
		\end{itemize}
		\item Each sense has a lexicographic category (\keyword{supersense}).
		\begin{itemize}
			\item E.g. \expword{05659525: noun.cognition}
		\end{itemize}
		\item WordNet represents semantic relations between senses.
	\end{itemize}
	
\end{frame}


\begin{frame}
	\frametitle{\insertshortsubtitle: \insertsection}
	\framesubtitle{\insertsubsection: Lexicographic categories (supersense)}
	
	\begin{block}{WordNet nouns lexicographic categories \cite{2019-jurafsky-martin} (Modified)}
		\fontsize{8}{12}\selectfont\bfseries
		\centering
		\begin{tblr}{
			colspec = {llllll},
			row{odd} = {lightblue},
			row{even} = {lightyellow},
			row{1} = {darkblue},
			rowsep=1pt,
			colsep=3pt,
		} 
			\textcolor{white}{Category} & \textcolor{white}{Example} & \textcolor{white}{Category} & \textcolor{white}{Example} &\textcolor{white}{Category} & \textcolor{white}{Example} \\
	%		\hline
			ACT & service & GROUP & team & PLANT & tree \\
			ANIMAL &  dog & LOCATION & area & POSSESSION & price \\
			ARTIFACT & car & MOTIVE & reason & PROCESS & process \\
			ATTRIBUTE & quality & NATURAL EVENT & earthquake & QUANTITY & amount \\
			BODY & hair & NATURAL OBJECT & flower & RELATION & connection \\
			COGNITION & idea & OTHER & stuff & SHAPE & square\\
			COMMUNICATION & review & PERSON & people & STATE & pain\\
			FEELING & discomfort & PHENOMENON & effect & SUBSTANCE & oil \\
			FOOD & bread & & & TIME & day\\
	%		\hline\hline
		\end{tblr}
	\end{block}
	
\end{frame}

\begin{frame}
	\frametitle{\insertshortsubtitle: \insertsection}
	\framesubtitle{\insertsubsection: Semantic relations (Nouns)}

	\begin{block}{Some nouns relations \cite{2019-jurafsky-martin} (Modified)}
		\fontsize{7}{14}\selectfont\bfseries\centering
		\begin{tblr}{
				colspec = {lll},
				row{odd} = {lightblue},
				row{even} = {lightyellow},
				row{1} = {darkblue},
				rowsep=1pt,
				colsep=4pt,
			} 
			\textcolor{white}{Relation} & \textcolor{white}{Definition} & \textcolor{white}{Example} \\
			Hypernym & from a specific concept to a more general one & \expword{breakfast\textsuperscript{1} \textrightarrow\ meal\textsuperscript{1} }\\
			Hyponym & from a general concept to a more specific one  & \expword{meal\textsuperscript{1} \textrightarrow\ lunch\textsuperscript{1}} \\
			Instance Hypernym & from an instance to its concept & \expword{Austen\textsuperscript{1} \textrightarrow\ author\textsuperscript{1}} \\
			Instance Hyponym & from a concept to its instances & \expword{composer\textsuperscript{1} \textrightarrow\ Bach\textsuperscript{1}} \\
			Part Meronym & from a whole to its part & \expword{table\textsuperscript{2} \textrightarrow\ leg\textsuperscript{3}} \\
			Part Holonym & from a part to its entirety & \expword{wheel\textsuperscript{7} \textrightarrow\ car\textsuperscript{1}} \\
			Antonym & between concepts with opposite meanings & \expword{leader\textsuperscript{1} $ \leftrightarrow $ follower\textsuperscript{1}}\\
			Derivation & from a word to a morphologically related word & \expword{destruction\textsuperscript{1} $ \leftrightarrow $ destroy\textsuperscript{1}} \\
		\end{tblr}
	\end{block}
	
\end{frame}

\begin{frame}
	\frametitle{\insertshortsubtitle: \insertsection}
	\framesubtitle{\insertsubsection: Semantic relations (Verbs)}
	
	\begin{block}{Some verbs relations \cite{2019-jurafsky-martin} (Modified)}
		\fontsize{7}{14}\selectfont\bfseries\centering
		\begin{tblr}{
				colspec = {lll},
				row{odd} = {lightblue},
				row{even} = {lightyellow},
				row{1} = {darkblue},
				rowsep=1pt,
				colsep=4pt,
			} 
			\textcolor{white}{Relation} & \textcolor{white}{Definition} & \textcolor{white}{Example} \\
			Hypernym & from a more specific verb to a more general one & \expword{fly\textsuperscript{9} \textrightarrow\ travel\textsuperscript{5}} \\
			Troponym & from a general verb to a more specific one & \expword{walk\textsuperscript{1} \textrightarrow\ stroll\textsuperscript{1}} \\
			Entails & from a verb to another event that is necessarily implied & \expword{snore\textsuperscript{1} \textrightarrow\ sleep\textsuperscript{1}} \\ 
			Antonym & between verbs with opposite meanings & \expword{increase\textsuperscript{1} $ \leftrightarrow $ decrease\textsuperscript{1}} \\
%			\hline\hline
		\end{tblr}
	\end{block}
	
\end{frame}

\begin{frame}
	\frametitle{\insertshortsubtitle: \insertsection}
	\framesubtitle{\insertsubsection: Some resources}
	
	\begin{itemize}
		\item WordNet APIs and libraries:
		\begin{itemize}
			\item NLTK (Python): \url{https://www.nltk.org/howto/wordnet.html}
			\item JWI (Java): \url{http://projects.csail.mit.edu/jwi}
			\item Wordnet (Ruby): \url{https://github.com/wordnet/wordnet}
			\item OpenNlp (.NET port, unofficial): \url{https://github.com/AlexPoint/OpenNlp}
		\end{itemize}
		\item Multilingual WordNet resources:
		\begin{itemize}
			\item Global WordNet Association: \url{http://globalwordnet.org/resources/wordnets-in-the-world/}
			\item Open Multilingual WordNet: \url{http://compling.hss.ntu.edu.sg/omw/}
		\end{itemize}
	\end{itemize}
	
\end{frame}

\begin{frame}
	\frametitle{\insertshortsubtitle: \insertsection}
	\framesubtitle{\insertsubsection: Similarity measures}
	
	\begin{itemize}
		\item Let us define some measures over concepts $c_i$ and $c_j$:
		\begin{itemize}
			\item $len(c_i, c_j)$: the length of the shortest path between $c_i$ and $c_j$ in the IS-A taxonomy;
			\item $lcs(c_i, c_j)$: the lowest common subsumer (most specific ancestor node);
			\item $depth(c_i)$: the depth of $c_i$ in the taxonomy (distance from the root);
			\item $depth_{max}$: the max $depth(c_i)$ of the taxonomy;
			\item $IC(c_i)$: the information content of $c_i$, defined as $-\log P(c_i)$ (Corpus-based).
		\end{itemize}
		\item Similarity measures:
		\begin{itemize}
			\item \optword{Path similarity}: $sim_{path}(c_i, c_j) = \frac{1}{1+len(c_i, c_j)}$
			\item \optword{Leacock Chodorow Similarity}: $sim_{lch}(c_i, c_j) = -\log \left( \frac{len(c_i, c_j)}{2 \cdot depth_{max}} \right)$
			\item \optword{Wu-Palmer Similarity}: $sim_{wup}(c_i, c_j) = \frac{2 \cdot depth(lcs(c_i, c_j))}{depth(c_i) + depth(c_j)}$
			\item \optword{Lin Similarity}: $sim_{lin}(c_i, c_j) = \frac{2*IC(lcs(c_i, c_j))}{IC(c_i)+IC(c_j)}$ 
			\item ...
		\end{itemize}
	\end{itemize}
	
\end{frame}

\begin{frame}
	\frametitle{\insertshortsubtitle: \insertsection}
	\framesubtitle{\insertsubsection: Exercise}
	
	\hgraphpage{wn-exp.pdf}
	
	\begin{itemize}
		\item Calculate path similarity between all concepts if applicable
	\end{itemize}
	
\end{frame}

\subsection{Other resources}

\begin{frame}
	\frametitle{\insertshortsubtitle: \insertsection}
	\framesubtitle{\insertsubsection: Other Lexical DBs}
	
	\begin{itemize}
		\item \optword{FrameNet} 
		\begin{itemize}
			\item It is based on frame semantics theory.
			\item A frame can be an event, a relation or an entity with its participants.
			\item E.g. \expword{The frame ``Cooking'' implies a cook (person), food being prepared, a container, and a heating instrument.}.
			\item Each frame is activated by a set of \keyword{lexical units}. E.g. \expword{bake, cook, fry, grill, prepare, roast}.
		\end{itemize}
		\item \optword{VerbNet}:
		\begin{itemize}
			\item It is a lexical base for verbs.
			\item It includes +30 principal thematic roles.
			\item The verbs are organized in classes.
		\end{itemize}
		\item \textbf{We will see these two DBs again in the next chapter.}
	\end{itemize}
	
\end{frame}

\begin{frame}
	\frametitle{\insertshortsubtitle: \insertsection}
	\framesubtitle{\insertsubsection: BabelNet}

	\begin{itemize}
		\item Multilingual sematic network
		\item \url{https://babelnet.org/}
	\end{itemize}
	
	\begin{figure}
		\hgraphpage{babelnet.pdf}
		\caption{BabelNet's structure \cite{2012-navigli-ponzetto}}
	\end{figure}
	
\end{frame}

%===================================================================================
\section{Vector representation of words}
%===================================================================================

\begin{frame}
	\frametitle{\insertshortsubtitle}
	\framesubtitle{\insertsection}

	\begin{itemize}
		\item A word can be represented using \keyword{One-Hot} encoding:
		\begin{itemize}
			\item One-Hot vectors do not encode semantic similarity (e.g., \expword{cat vs. dog}) or contextual proximity (e.g., \expword{coffee vs. cup}).
			\item How to represent a document/sentence using this representation?
		\end{itemize}
	
		\item \optword{Term-document}:
		\begin{itemize}
			\item A term is represented using the documents containing it (or the inverse).
			\item E.g. \expword{Term frequency - inverse document frequency (Tf-Idf)}.
		\end{itemize}
	
		\item \optword{Term-term}:
		\begin{itemize}
			\item A term is represented using co-occurrence with other terms.
		\end{itemize}
	
		\item \optword{Term-concept-document}:
		\begin{itemize}
			\item Terms and documents are represented using a vector of concepts.
			\item E.g. \expword{Latent Semantic Analysis (LSA)}.
			\item Welcome to \keyword{Word embedding}!
			\begin{itemize}
				\item Word embeddings are dense vector representations of words, not necessarily learned via neural networks.
			\end{itemize}
		\end{itemize}
	
	\end{itemize}

\end{frame}


\subsection{TF-IDF}

\begin{frame}
	\frametitle{\insertshortsubtitle: \insertsection}
	\framesubtitle{\insertsubsection}

	\begin{itemize}
		\item A document can be represented using the words it contains.
		\item Word's frequency in a document is called \keyword{term frequency (TF)}.
		\item There are words that are repeated a lot, which have no added meaning (e.g. \expword{
		Prepositions}).
		\item Words that appear in multiple documents are less important.
		\item High TF-IDF → word is frequent in this document but rare in the corpus → likely important.
	\end{itemize}
	
	\begin{block}{TF-IDF calculation}
		\[
		TF_d(t) =  |\{t_i \in d / t_i = t\}|
		\hskip2cm 
		DF_D(t) = |\{d \in D / t \in d\}|
		\]
		\[IDF_D(t) = \log_{10} \left( \frac{|\{d \in D\}|}{DF_D(t)} \right)\]
		\[TF\text{-}IDF_{d, D}(t) = TF_d(t) * IDF_D(t)\]
	\end{block}

\end{frame}


\begin{frame}
	\frametitle{\insertshortsubtitle: \insertsection}
	\framesubtitle{\insertsubsection: Example of TF representation}

	\begin{exampleblock}{Example of some sentences}
		\begin{itemize}
			\item S1: a computer can help you
			\item S2: he can help you and he wants to help you
			\item S3: he wants a computer and a computer for you
		\end{itemize}
	\end{exampleblock}
	
	\begin{center}
		\begin{tabular}{lllllllllll}
		\hline\hline
		& a & and & can & computer & for & he & help & to & wants & you \\
		\hline
		S1 & 1 & 0 & 1 & 1 & 0 & 0 & 1 & 0 & 0 & 1\\
		S2 & 0 & 1 & 1 & 0 & 0 & 2 & 2 & 1 & 1 & 2\\
		S3 & 2 & 1 & 0 & 2 & 1 & 1 & 0 & 0 & 1 & 1\\
		\hline\hline
	\end{tabular}
	\end{center}

\end{frame}


\begin{frame}
	\frametitle{\insertshortsubtitle: \insertsection}
	\framesubtitle{\insertsubsection: Cosine similarity}

	\begin{minipage}{.68\textwidth}
	\begin{itemize}
		\item Cosine similarity measures the angle between the vectors; 1 means identical orientation, 0 means orthogonal (no similarity).
		\item Let $a$ and $b$ be two documents represented by two vectors $\overrightarrow{a}$ and $\overrightarrow{b}$ respectively.
		\item The representation can be vocabulary words (TF or TF-IDF) or other document features.
		\item Let the vectors have length $n$.
	\end{itemize}
	\end{minipage}
	\begin{minipage}{.3\textwidth}
	\hgraphpage{exp-cos.pdf}
	\end{minipage}
	
	\begin{block}{Cosine similarity between two vectors}
		\[
		Cos(\theta) = \frac{\overrightarrow{a} \overrightarrow{b}}{||\overrightarrow{a}||\, ||\overrightarrow{b}||}
		= \frac{\sum_{i=1}^{n} a_i b_i}{\sqrt{\sum_{i=1}^{n} a_i^2} \sqrt{\sum_{i=1}^{n} b_i^2}}
		\]
	\end{block}

\end{frame}

\subsection{Term-term}

\begin{frame}
	\frametitle{\insertshortsubtitle: \insertsection}
	\framesubtitle{\insertsubsection}

	\begin{itemize}
		\item Words can be represented using other vocabulary words based on co-occurrence.
		\item To represent all words of a vocabulary $ V $, a matrix $|V| \times |V|$ must be used.
		\item Each word is represented by a vector of length $|V|$, where each entry counts its co-occurrence with a \keyword{context} word.
		\item Co-occurrence can be calculated according to documents, sentences or windows around the word.
		\item For example, \expword{a context window of size 2-2 considers the 2 words before and 2 words after the target word}.
	\end{itemize}

\end{frame}

\begin{frame}
	\frametitle{\insertshortsubtitle: \insertsection}
	\framesubtitle{\insertsubsection: Example of a window 2-2}

	\begin{exampleblock}{Example of some sentences}
		\begin{itemize}
			\item S1: a computer can help you
			\item S2: he can help you and he wants to help you
			\item S3: he wants a computer and a computer for you
		\end{itemize}
	\end{exampleblock}
	
	\begin{center}
		\scriptsize
		\begin{tabular}{lllllllllll}
			\hline\hline
			& a & and & can & computer & for & he & help & to & wants & you \\
			\hline
			a & 0 & 2 & 1 & 4 & 1 & 1 & 0 & 0 & 1 & 0\\
			and & 2 & 0 & 0 & 2 & 0 & 1 & 1 & 0 & 1 & 1\\
			can & 1 & 0 & 0 & 1 & 0 & 1 & 2 & 0 & 0 & 2\\
			computer & 4 & 2 & 1 & 0 & 1 & 0 & 1 & 0 & 1 & 1\\
			for & 1 & 0 & 0 & 1 & 0 & 0 & 0 & 0 & 0 & 1\\
			he & 1 & 1 & 1 & 0 & 0 & 0 & 1 & 1 & 2 & 1\\
			help & 0 & 1 & 2 & 1 & 0 & 1 & 0 & 1 & 1 & 3\\
			to & 0 & 0 & 0 & 0 & 0 & 1 & 1 & 0 & 1 & 1\\
			wants & 1 & 1 & 0 & 1 & 0 & 2 & 1 & 1 & 0 & 0\\
			you & 0 & 1 & 2 & 1 & 1 & 1 & 3 & 1 & 0 & 0\\
			\hline\hline
		\end{tabular}
	\end{center}

\end{frame}

\begin{frame}
	\frametitle{\insertshortsubtitle: \insertsection}
	\framesubtitle{\insertsubsection: Cosine similarity}

	\begin{itemize}
		\item Similarity between two words can be calculated using cosine similarity (previously seen) or other similarity measures such as: Euclidean distance and dot product.
	\end{itemize}
	
	\begin{exampleblock}{Example of cosine similarity between ``computer'' and ``you''}
		\fontsize{6}{16}\selectfont\bfseries\boldmath
		\begin{align*}
		& \text{Cos(``computer'', ``you'')} \\ 
		& = \frac{4*0 + 2*1 + 1*2 + 0*1 + 1*1 + 0*1 + 1*3 + 0*1 + 1*0 + 1*0}{\sqrt{4^2 + 2^2 + 1^2 + 0^2 + 1^2 + 0^2 + 1^2 + 0^2 + 1^2 + 1^2} \sqrt{0^2 + 1^2 + 2^2 + 1^2 + 1^2 + 1^2 + 3^2 + 1^2 + 0^2 + 0^2}}\\
		& = \frac{8}{\sqrt{25}\sqrt{18}}\\
		& = \frac{8}{3*5*\sqrt{2}}\\
		\end{align*}
	\end{exampleblock}

\end{frame}

\subsection{Latent semantic analysis (LSA)}

\begin{frame}
	\frametitle{\insertshortsubtitle: \insertsection}
	\framesubtitle{\insertsubsection}

	\begin{itemize}
		\item The term-document matrix $X \in \mathbb{R}^{N \times M}$ is usually:
		\begin{itemize}
			\item high-dimensional;
			\item sparse (many zeros, see TF example).
		\end{itemize}
		\item $L$ latent (hidden) concepts (topics) can be used to represent:
		\begin{itemize}
			\item terms: $T \in \mathbb{R}^{N \times L}$;
			\item documents: $D \in \mathbb{R}^{M \times L}$.
		\end{itemize}
		\item $L \le \min(N, M)$
		\item $X$ can be decomposed using \keyword{Singular value decomposition (SVD)}
		\begin{itemize}
			\item $X = T \times S \times D^\top$, where $S$ is diagonal with singular values.
		\end{itemize}
	\end{itemize}

\end{frame}


\begin{frame}
	\frametitle{\insertshortsubtitle: \insertsection}
	\framesubtitle{\insertsubsection: SVD formulation}
	
%	\[
%	X_{N \times M} = 
%	\overbrace{T_{N \times L}}^{\text{term-concept}}
%	\; 
%	\overbrace{S_{L \times L}}^{\text{singular values}}
%	\;
%	\overbrace{D_{M \times L}^\top}^{\text{concept-document}}
%	\]
	
	\begin{itemize}
		\item $T^\top T = \mathbb{I}_{L \times L}$, $D^\top D = \mathbb{I}_{L \times L}$
		\item $S$ is diagonal with decreasing values: $s_{11} \ge s_{22} \ge ... \ge s_{LL} > 0$
		\item Computation methods: \keyword{Lanczos algorithm}, \keyword{QR algorithm}
	\end{itemize}
	
	
	\begin{block}{SVD of term-document matrix}
		\scriptsize\bfseries
		\[
		\overbrace{
			\begin{bmatrix}
			x_{11} & \ldots & \ldots & \ldots & x_{1M} \\ 
			\vdots & \ddots & \ddots & \ddots &\vdots \\
			\vdots & \ddots & \ddots & \ddots &\vdots \\
			x_{N1} & \ldots & \ldots & \ldots & x_{NM} \\ 
			\end{bmatrix}
		}^{X \text{: term-document}}
		=
		\overbrace{
			\left[
			\begin{bmatrix}
			t_{11} \\ 
			\vdots \\
			\vdots \\
			t_{N1} \\ 
			\end{bmatrix}
			\begin{matrix}
			\ldots \\ 
			\end{matrix}
			\begin{bmatrix}
			t_{1L} \\ 
			\vdots \\
			\vdots \\
			t_{NL} \\ 
			\end{bmatrix}
			\right]
		}^{T \text{: term-concept}}
		\times 
		\overbrace{
			\begin{bmatrix}
			s_{11} & \ldots & 0 \\
			0 & \ddots & 0 \\
			0 & \ldots & s_{LL} \\
			\end{bmatrix}
		}^{S \text{: concept-concept}}
		\times 
		\overbrace{
			\begin{bmatrix}
			\begin{bmatrix}
			d_{11} & \ldots & \ldots & \ldots & d_{1M} \\
			\end{bmatrix}\\
			\vdots \\
			\begin{bmatrix}
			d_{L1} & \ldots & \ldots & \ldots & d_{LM} \\
			\end{bmatrix}\\
			\end{bmatrix}
		}^{D^\top \text{: concept-document}}
		\]
		
	\end{block}
	
\end{frame}


%===================================================================================
\section{Word embedding}
%===================================================================================

\begin{frame}
	\frametitle{\insertshortsubtitle}
	\framesubtitle{\insertsection}
	
	\begin{itemize}
		\item Document-term and term-term matrices are large and sparse, making them computationally expensive to store and process.
		\item Using \keyword{Latent Semantic Analysis (LSA)}, words can be represented by low-dimensional vectors in a continuous space.
		\item This representation is referred to as a \keyword{word embedding}.
		\item In this course, we focus on neural network-based word embeddings.
		\item Such embeddings capture semantic and syntactic relations between words, including similarity and analogy.
	\end{itemize}
	
\end{frame}


\begin{frame}
	\frametitle{\insertshortsubtitle}
	\framesubtitle{\insertsection: Some humor}
	
	\begin{center}
			\vgraphpage{humor/humor-embeddings.jpeg}	
	\end{center}

\end{frame}


\subsection{Word2vec}

\begin{frame}
	\frametitle{\insertshortsubtitle: \insertsection}
	\framesubtitle{\insertsubsection}
	
	\begin{itemize}
		\item In term-term representations, a word is represented based on its neighboring words (\keyword{context}):
		\begin{itemize}
			\item Most words co-occur with only a small subset of the vocabulary, resulting in sparse vectors;
			\item The vectors are high-dimensional and mostly zero-valued.
		\end{itemize}
		\item Word2Vec (Google) implements two neural network-based methods for \keyword{word embedding} \cite{2013-mikolov-al}:
		\begin{itemize}
			\item \optword{CBOW}: Continuous Bag-of-Words, which predicts a word from its context.
			\item \optword{Skip-gram}: predicts context words from a target word.
		\end{itemize}
		\item Each word is encoded as a dense, low-dimensional vector (50–1000 dimensions) that captures semantic similarity.
	\end{itemize}
	

\end{frame}

\begin{frame}
	\frametitle{\insertshortsubtitle: \insertsection}
	\framesubtitle{\insertsubsection: CBOW}
	
	\begin{minipage}{.63\textwidth}
		
		\begin{itemize}
			\item Predict the target word $w_i$ from its surrounding context.
			\item The context window consists of the two preceding and two succeeding words.
			\item Let $T$ denote the total number of words in the corpus.
		\end{itemize}
		
		\begin{block}{CBOW Loss Function}
			\[
			J(\theta) = -\frac{1}{T} \sum_{i=1}^{T} \log p(w_i \mid w_{i-2}, w_{i-1}, w_{i+1}, w_{i+2})
			\]
		\end{block}
		
	\end{minipage}
	\begin{minipage}{.36\textwidth}
		\hgraphpage{word2vec-cbow.pdf}
	\end{minipage}
	
\end{frame}

\begin{frame}
	\frametitle{\insertshortsubtitle: \insertsection}
	\framesubtitle{\insertsubsection: Skip-gram}
	
	\begin{minipage}{.58\textwidth}
		\begin{itemize}
			\item Predict the surrounding context words given the target word $w_i$.
			\item The context window consists of the two preceding and two succeeding words.
			\item Let $T$ denote the total number of words in the corpus.
		\end{itemize}
		
		\begin{block}{Skip-gram Loss Function}
			\[
			J(\theta) = -\frac{1}{T} \sum_{i=1}^{T} \sum_{j=i-2, j \ne i}^{i+2} \log p(w_j \mid w_i)
			\]
		\end{block}
		
	\end{minipage}
	\begin{minipage}{.4\textwidth}
		\hgraphpage{word2vec-skip.pdf}
	\end{minipage}
	
\end{frame}

\subsection{GloVe}

\begin{frame}
	\frametitle{\insertshortsubtitle: \insertsection}
	\framesubtitle{\insertsubsection}
	
	\begin{itemize}
		\item \keyword{GloVe}: Global Vectors, developed at Stanford \cite{2014-pennington-al}.
		\item GloVe leverages both approaches: global co-occurrence statistics (as in LSA) and local context information (as in CBOW).
		\item Let $X[V, V]$ be the term-term co-occurrence matrix, where:
		\begin{itemize}
			\item $X_{ij}$ = number of times word $w_j$ appears in the context of word $w_i$;
			\item $X_i = \sum_j X_{ij}$ is the total number of occurrences of $w_i$ in the corpus.
		\end{itemize}
		\item The probability of observing $w_j$ in the context of $w_i$ is:
		\[
		P_{ij} = \frac{X_{ij}}{X_i}
		\]
	\end{itemize}
	

\end{frame}

\begin{frame}
	\frametitle{\insertshortsubtitle: \insertsection}
	\framesubtitle{\insertsubsection: Motivation}
	
	\vspace{6pt}
	\begin{minipage}{.5\textwidth}
		\begin{itemize}
			\item Consider $w_i = \text{ice}$ and $w_j = \text{steam}$.
			\item Let $p_{ik}, p_{jk}$ denote the probabilities of context word $w_k$ occurring with $w_i$ and $w_j$ respectively.
			\item Define the co-occurrence ratio $\frac{p_{ik}}{p_{jk}}$.
		\end{itemize}
		
	\end{minipage}
	\begin{minipage}{.48\textwidth}
			\fontsize{6}{10}\selectfont\bfseries
			\begin{tblr}{
					colspec = {p{1.2cm}@{\hskip3pt}p{1cm}@{\hskip3pt}p{1cm}@{\hskip3pt}p{1cm}@{\hskip3pt}p{1.2cm}@{\hskip0pt}},
					row{odd} = {lightblue},
					row{even} = {lightyellow},
					row{1} = {darkblue},
					colsep = 0pt,
					rowsep= 1pt,
				} 
	%			\hline\hline
				
				\textcolor{white}{\textbf{Probability and Ratio}} & \textcolor{white}{\textbf{k = solid}} & \textcolor{white}{\textbf{k = gas}} & \textcolor{white}{\textbf{k = water}} & \textcolor{white}{\textbf{k = fashion}} \\
	%			\hline
				P(k\textbar ice) & 1.9 * 10\textsuperscript{-4} & 6.6 * 10\textsuperscript{-5} & 3.0 * 10\textsuperscript{-3} & 1.7 * 10\textsuperscript{-5} \\
				P(k\textbar steam) & 2.2 * 10\textsuperscript{-5} & 7.8 * 10\textsuperscript{-4} & 2.2 * 10\textsuperscript{-3} & 1.8 * 10\textsuperscript{-5} \\
	%			\hline
				\boldmath\footnotesize$\frac{P(k|ice)}{P(k|steam)}$ & 8.9 & 8.5 * 10\textsuperscript{-2} & 1.36 & 0.96 \\
	%			\hline\hline
			\end{tblr}
	%		\caption{Exemple de probabilités et un ratio \cite{2014-pennington-al}}
	%	\end{figure}
	\end{minipage}
		
	\begin{itemize}
		\item Ratio's value follows relations $R(w_k, w_i)$ and $R(w_k, w_j)$:
		\begin{itemize}
			\item If $R(w_k, w_i) \wedge \neg R(w_k, w_j)$, the ratio will be large. E.g. \expword{solid}.
			\item If $\neg R(w_k, w_i) \wedge R(w_k, w_j)$, the ratio will be small. E.g. \expword{gas}.
			\item If $R(w_k, w_i) \wedge R(w_k, w_j)$, the ratio will be close to $1$. E.g. \expword{water}.
			\item If $\neg R(w_k, w_i) \wedge \neg R(w_k, w_j)$, the ratio will be close to $1$. E.g. \expword{fashion}.
		\end{itemize}
		\item The idea is to train a function $F$ on $w_i$, $w_j$ and $\tilde{w_k}$ representations:
		\[F(w_i, w_j, \tilde{w_k}) = \frac{p_{ik}}{p_{jk}}\]
	\end{itemize}
	
\end{frame}

\begin{frame}
	\frametitle{\insertshortsubtitle: \insertsection}
	\framesubtitle{\insertsubsection: Formulation (intuition)}
	
	\begin{itemize}
		\item We impose a constraint on $F$ using the difference of word vectors:
		\[
		F(w_i - w_j, \tilde{w_k}) = \frac{p_{ik}}{p_{jk}}
		\]
		\item Map the vector difference to a scalar via the dot product:
		\[
		F((w_i - w_j)^\top \tilde{w_k}) = \frac{p_{ik}}{p_{jk}}
		\]
		\item Considering symmetry ($w \leftrightarrow \tilde{w}$ and $X \leftrightarrow X^\top$), 
		$F$ must satisfy the homomorphism property:
		\[
		F((w_i - w_j)^\top \tilde{w_k}) = \frac{F(w_i^\top \tilde{w_k})}{F(w_j^\top \tilde{w_k})}
		\]
		\item Therefore, we set:
		$F(w_i^\top \tilde{w_k}) = p_{ik} = \frac{X_{ik}}{X_i}$
	
		\item Choosing $F = \exp$ satisfies the property, yielding:
		\[w_i^\top \tilde{w_k} = \log X_{ik} - \log X_i\]
		
	\end{itemize}
	
	
\end{frame}

\begin{frame}
	\frametitle{\insertshortsubtitle: \insertsection}
	\framesubtitle{\insertsubsection: Formulation (Final version)}
	
	\begin{itemize}
		\item $\log X_i$ is independent of $j$, so we introduce a bias term $b_i$ for $w_i$.
		\item To maintain symmetry, we also introduce a bias $\tilde{b_j}$ for $\tilde{w_j}$.
		\item The loss function $J$ is based on \keyword{weighted least squares}:
		\begin{itemize}
			\item Co-occurrences contribute unequally: rare co-occurrences have less impact.
			\item We apply a weighting function $f(X_{ij})$ to control the contribution.
		\end{itemize}
	\end{itemize}
	
	\vskip-2pt\begin{block}{GloVe Formulation}\vskip-2pt
		\[w_i^\top \tilde{w_j} + b_i + \tilde{b_j} = \log X_{ij} \]
		\[J(\theta) = \sum_{i=1}^{V} \sum_{j=1}^{V} f(X_{ij}) (w_i^\top \tilde{w_j} + b_i + \tilde{b_j} - \log X_{ij})^2\]
		\vskip-2pt\[f(x) = \begin{cases}
			\frac{x}{x_{max}} & \text{if } x < x_{max} \\
			1 & \text{otherwise}
		\end{cases}\]
	\end{block}
	
\end{frame}

\begin{frame}
	\frametitle{\insertshortsubtitle: \insertsection}
	\framesubtitle{\insertsubsection: Architecture}
	
	\begin{center}
		\vgraphpage{glove.pdf}
	\end{center}
	
\end{frame}

\subsection{Contextual embedding}

\begin{frame}
	\frametitle{\insertshortsubtitle: \insertsection}
	\framesubtitle{\insertsubsection: Contextual embedding}
	
	\begin{exampleblock}{Example of polysemy (word with multiple senses)}
		\begin{itemize}
			\item My cat is always chasing a \expword{mouse}. (rodent)
			\item The computer's \expword{mouse} is a pointing device. (electronic device)
			\item You are such a \expword{mouse}! (timid person)
		\end{itemize}
	\end{exampleblock}
	
	\begin{itemize}
		\item To what extent can Word2vec or GloVe capture the semantic relations between \expword{mouse} and words such as \expword{cat}, \expword{computer}, or \expword{timid}? Explain why.
		\item Are these embeddings capable of distinguishing between the different senses of \expword{mouse}? Justify your answer.
	\end{itemize}
	
\end{frame}

\begin{frame}
	\frametitle{\insertshortsubtitle: \insertsection}
	\framesubtitle{\insertsubsection: ELMo}
	
	\begin{minipage}{.65\textwidth}
		\begin{itemize}
			\item \keyword{ELMo}: Embeddings from Language Models
			\item \url{http://allennlp.org/elmo} 
			\item \cite{2018-peters-al}
			\item Key Characteristics of ELMo Representations:
			\begin{itemize}
				\item \textbf{Contextual}: captures the meaning of a word based on its full sentence context.
				\item \textbf{Deep}: integrates all intermediate representations from the neural network layers.
				\item \textbf{Character-based}: captures morphological features and handles out-of-vocabulary words.
			\end{itemize}
		\end{itemize}
	\end{minipage}
	\begin{minipage}{.33\textwidth}
		\vspace{1.8cm}
		\hgraphpage{Elmo-img.jpg}
	\end{minipage}
	
\end{frame}



\begin{frame}
	\frametitle{\insertshortsubtitle: \insertsection}
	\framesubtitle{\insertsubsection: ELMo (Architecture)}
	\vspace{-4pt}
%	\begin{center}
		\hgraphpage{elmo-arch.pdf}
%	\end{center}
	
\end{frame}

\begin{frame}
	\frametitle{\insertshortsubtitle: \insertsection}
	\framesubtitle{\insertsubsection: ELMo (Description)}
%	\small
	\begin{itemize}
		\item Each word $w_k$ is represented by an embedding $x_k$, computed from its character-level representation (\cite{2015-kim-al}).
		\item The model uses $L$ layers of bidirectional LSTMs to capture context.
		\item For a sentence of length $N$, the model maximizes:
		\begin{center}
			$\sum_{k=1}^{N} (
			\log P(w_k | w_1,\ldots,w_{k-1}; \Theta_x, \overrightarrow{\Theta}_{LSTM}, \Theta_s)
			+
			\log P(w_k | w_{k+1},\ldots,w_{N}; \Theta_x, \overleftarrow{\Theta}_{LSTM}, \Theta_s)
			)
			$
		\end{center}
		
		\item The set of representations for $w_k$ is:
		\[
		R_k = \{x_k^{LM}, \overrightarrow{h}_{LM}^{k, j}, \overleftarrow{h}_{LM}^{k, j} \mid j= 1 \ldots L \} 
		= \{ h_{LM}^{k, j} \mid j=0 \ldots L \}
		\]
		
		\item To adapt ELMo for a downstream task, additional parameters $\Theta^{task}$ are trained to combine the intermediate representations into a single task-specific embedding:
		\vspace{-8pt}
		\[
		ELMo_k^{task} = E(R_k; \Theta^{task}) = \gamma^{task} \sum_{j=0}^{L} \theta_j^{task} h_{LM}^{k, j}
		\]
	\end{itemize}
	
\end{frame}


\begin{frame}
	\frametitle{\insertshortsubtitle: \insertsection}
	\framesubtitle{\insertsubsection: ELMo (Some humor)}

	\begin{center}
		\vgraphpage{humor/humor-elmo.png}
	\end{center}
	
\end{frame}


\begin{frame}
	\frametitle{\insertshortsubtitle: \insertsection}
	\framesubtitle{\insertsubsection: BERT}
	
	\begin{minipage}{.71\textwidth}
		\begin{itemize}
			\item \keyword{BERT}: Bidirectional Encoder Representations from Transformers
			\item \url{https://github.com/google-research/bert}
			\item \cite{2019-devlin-al}
			\item Representation characteristics:
			\begin{itemize}
				\item \textbf{Contextual}: embeddings depend on the entire sentence;
				\item \textbf{Fully bidirectional}: captures relationships among all words using self-attention;
				\item \textbf{Token-based}: uses WordPiece subword tokenization to handle rare words and morphological variations;
				\item \textbf{Supports transfer learning}: pre-trained on large corpora, fine-tuned for downstream tasks.
			\end{itemize}
		\end{itemize}
	\end{minipage}
	\begin{minipage}{.27\textwidth}
		\vspace{1.1cm}
		\hgraphpage{Bert-img.png}
	\end{minipage}
	
\end{frame}


\begin{frame}
	\frametitle{\insertshortsubtitle: \insertsection}
	\framesubtitle{\insertsubsection: BERT (Architecture)}
	
	\begin{center}
		\vgraphpage{bert-arch.pdf}
	\end{center}
	
\end{frame}

\begin{frame}
	\frametitle{\insertshortsubtitle: \insertsection}
	\framesubtitle{\insertsubsection: BERT (Input)}
	
	\begin{itemize}
		\item Text is tokenized using \keyword{WordPiece} \cite{2016-wu-al}.
		\item Input sequence of tokens has maximum length $T = 512$:
		\begin{itemize}
			\item The first token is a special token ``\keyword{[CLS]}'' for classification.
			\item The remaining tokens correspond to the sentence(s), with ``\keyword{[SEP]}'' separating sentences or segments.
		\end{itemize}
		\item Each token is represented by a vector of size $N$, obtained by summing three embeddings:
		\begin{itemize}
			\item \optword{Token embedding}: learned embedding for each vocabulary token (size $V \times N$).
			\item \optword{Position embedding}: learned embedding for each position in the sequence (size $T \times N$).
			\item \optword{Segment embedding}: learned embedding for sentence segment (size $2 \times N$ for sentence1 or sentence2).
		\end{itemize}
	\end{itemize}
	
\end{frame}


\begin{frame}
	\frametitle{\insertshortsubtitle: \insertsection}
	\framesubtitle{\insertsubsection: BERT (Pre-training)}
	
	\begin{itemize}
		\item The model is based on \keyword{Transformers} (encoder only) \cite{2017-vaswani-al}.
		\item Pre-trained on two tasks:
		\begin{itemize}
			\item \optword{Masked Language Model (MLM)}:
			\begin{itemize}
				\item Randomly mask 15\% of tokens in input.
				\item Use special token \keyword{[MASK]} (appears only in pre-training).
				\item 80\% replaced by \keyword{[MASK]}, 10\% by a random token, 10\% left unchanged.
			\end{itemize}
			\item \optword{Next Sentence Prediction (NSP)}:
			\begin{itemize}
				\item Given two input sentences, predict whether the second follows the first.
				\item For classification, the \keyword{[CLS]} token's final hidden state is used to predict ``\textbf{IsNext}'' or ``\textbf{NotNext}''.
			\end{itemize}
		\end{itemize}
		\item After pre-training, BERT can be fine-tuned for downstream tasks
	\end{itemize}
	
\end{frame}


\begin{frame}
	\frametitle{\insertshortsubtitle: \insertsection}
	\framesubtitle{\insertsubsection: BERT (Fine-tuning)}

	\begin{center}
		\vgraphpage{bert-tasks.pdf}
	\end{center}
	
\end{frame}

\begin{frame}
	\frametitle{\insertshortsubtitle: \insertsection}
	\framesubtitle{\insertsubsection: BERT (Some humor)}

	\begin{center}
		\vgraphpage{humor/humor-bert.jpg}
	\end{center}
	
\end{frame}

\begin{frame}
	\frametitle{\insertshortsubtitle: \insertsection}
	\framesubtitle{\insertsubsection: Intrinsic evaluation}
	
	\begin{itemize}
		\item \optword{WordSimilarity-353 Test Collection} \cite{2002-finkelstein-al}:
		\begin{itemize}
			\item Manually annotated word pairs with similarity scores from 0 to 10.
			\item Model performance is measured using Spearman correlation between cosine similarities and human judgments.
		\end{itemize}
		
		\item \optword{SimLex-999} \cite{2015-hill-al}:
		\begin{itemize}
			\item Similar to WordSimilarity-353 but focused on genuine similarity rather than association.
			\item E.g. \expword{coast, shore} are similar, \expword{clothes, closet} are associated but not similar.
		\end{itemize}
		
		\item \optword{Word analogies} \cite{2013-mikolov-al2}:
		\begin{itemize}
			\item Dataset format: $(w_{i1}:w_{j1} :: w_{i2}:w_{j2})$
			\item Tests the ability of embeddings to capture semantic relations: $w_{j2} \approx w_{i1} - w_{i2} + w_{j1}$.
			\item E.g. \expword{(King:Queen :: Man:Woman) $\Rightarrow$ King - Man + Woman $\approx$ Queen}
		\end{itemize}
		
	\end{itemize}
	
\end{frame}


\begin{frame}
	\frametitle{\insertshortsubtitle: \insertsection}
	\framesubtitle{\insertsubsection: Extrinsic evaluation and bias}
	
	\begin{itemize}
		\item \optword{Extrinsic evaluation}: Assess models based on their performance on downstream tasks.
		\begin{itemize}
			\item Compare a new model with a baseline on a specific task.
			\item E.g. GloVe outperformed LSA on Named Entity Recognition (NER) \cite{2014-pennington-al}.
		\end{itemize}
		\item \optword{Bias in embeddings}: Word embeddings can encode and reproduce societal stereotypes.
		\begin{itemize}
			\item Gendered associations: \expword{she → homemaker, nurse, receptionist}, \expword{he → maestro, skipper, protege} \cite{2017-caliskan-al}.
			\item Such biases can negatively impact the performance of downstream NLP tasks.
			\item E.g. \expword{Anaphora resolution may fail to correctly link the pronoun ``she'' with a professional role such as ``doctor''}.
		\end{itemize}
	\end{itemize}
	
\end{frame}


%===================================================================================
\section{Word Sense Disambiguation (WSD)}
%===================================================================================

\begin{frame}
	\frametitle{\insertshortsubtitle}
	\framesubtitle{\insertsection}
	
	\begin{itemize}
		\item Determining the intended sense of a word in context.
		\item WSD is important for several NLP tasks:
		\begin{itemize}
			\item \textbf{Parsing:}
			\begin{itemize}
				\item \expword{I \underline{fish} in the river} (Is ``fish'' a verb or a noun?)
				\item \expword{The \underline{fish} was too big} (Is ``fish'' a verb or a noun?)
			\end{itemize}
			\item \textbf{Machine translation:}
			\begin{itemize}
				\item \expword{I withdrew money from the \underline{bank}} (Does ``bank'' refer to a financial institution or a river edge?)
				\item \expword{I fish on the \underline{bank}} (Does ``bank'' refer to a financial institution or a river edge?)
			\end{itemize}
		\end{itemize}
	\end{itemize}
	
\end{frame}

\subsection{Knowledge bases based}

\begin{frame}
	\frametitle{\insertshortsubtitle: \insertsection}
	\framesubtitle{\insertsubsection: Lesk algorithm}
	
	\begin{block}{Lesk algorithm}
		\footnotesize\vspace{-3pt}
		\begin{algorithm}[H]
			\KwData{a word $w$; a sentence $s$ containing $w$}
			\KwResult{The sense of $w$ in $s$}
			
			best\_sense \textleftarrow most frequent sense of $w$\;
			overlap\_max \textleftarrow 0\;
			context \textleftarrow set of words in $s$\; 
			
			\ForEach{sense $w_i$ of $w$}{ 
				signature \textleftarrow set of words in the gloss and example usages of sense $w_i$\;
				overlap \textleftarrow number of common words between \textbf{context} and \textbf{signature}\;
				\If{overlap $>$ overlap\_max}{
					overlap\_max \textleftarrow overlap\;
					best\_sense \textleftarrow $w_i$\;
				}
			}
			
			\Return best\_sense \;
		\end{algorithm}
	\end{block}
\end{frame}

\begin{frame}
	\frametitle{\insertshortsubtitle: \insertsection}
	\framesubtitle{\insertsubsection: Graph-based}
	
	\begin{itemize}
		\item Example: \url{http://babelfy.org/} \cite{2014-moro-al}
		\item \optword{Step 1: Semantic signature construction}
		\begin{itemize}
			\item Assign a weight to each edge connecting two concepts, proportional to the number of triangles they participate in.
			\item Compute the conditional probability of a concept given a neighboring concept using these weights.
			\item Propagate importance using the \keyword{Random Walk with Restart} algorithm.
		\end{itemize}
		\item \optword{Step 2: Candidate identification}
		\begin{itemize}
			\item Apply \keyword{POS tagging} to the input text.
			\item Extract all possible senses of words and multi-word expressions.
		\end{itemize}
		\item \optword{Step 3: Candidate disambiguation}
		\begin{itemize}
			\item Construct a graph based on semantic signatures and candidate senses.
			\item Build a subgraph by pruning edges with low weights.
		\end{itemize}
	\end{itemize}
	
\end{frame}

\subsection{ML-based}

\begin{frame}
	\frametitle{\insertshortsubtitle: \insertsection}
	\framesubtitle{\insertsubsection: Supervised ML}
	
	\begin{itemize}
		\item Uses an annotated corpus (e.g., \expword{SemCor}).
		\item Each word is annotated with the corresponding sense ID from a lexical resource (e.g., \expword{WordNet}).
		\item Example: \expword{You will find9 that avocado1 is1 unlike1 other1 fruit1 you have ever1 tasted2}.
		\item The model outputs a One-Hot encoded vector over all possible senses.
		\item Features can include the target word as well as its surrounding context words.
		\item Disambiguating all words in a text can be approached as a sequence labeling task, similar to POS tagging.
	\end{itemize}
	
\end{frame}

\begin{frame}
	\frametitle{\insertshortsubtitle: \insertsection}
	\framesubtitle{\insertsubsection: Contextual embedding}
	
	\begin{minipage}{.68\textwidth}
		\begin{itemize}
			\item 1-nearest-neighbor algorithm
			\item Using contextual embeddings (\expword{ELMo, BERT, etc.})
			\item Train a model on an annotated corpus to obtain embeddings for each token.
			\item To obtain a sense embedding $v_s$, compute the average of all token embeddings $c_i$ corresponding to that sense:
			\[
			v_s = \frac{1}{n} \sum_{i=1}^{n} c_i
			\] 
		\end{itemize}
	\end{minipage}
	\begin{minipage}{.3\textwidth}
		\begin{figure}
			\hgraphpage{exp-wsd-nn.pdf}
			\caption{Example of WSD with 1-nearest-neighbor \cite{2019-jurafsky-martin}}
		\end{figure}
	\end{minipage}
	
\end{frame}


\insertbibliography{NLP07}{*}

\end{document}

